% 原书第 19—20 页；第 2.2 节及三幅因果图。
\section{因果图}
\label{sec:ch02-causal-graphs}

上一节将 SCM 介绍为一种系统的方法：从随机输入出发，为一组变量赋值，界定所需的因果假设，并评价因果目标参数的可识别性（identifiability）。因果图（causal graph）则是另一种常用表示方法，可以表达 SCM 所编码的部分假设。本书所有的因果图都是有向无环图（directed acyclic graph，DAG）：连接节点的边仅有一个箭头，而且不存在闭合回路。然而，面对纵向数据（longitudinal data）、网络及其他复杂数据结构时，这种表示方式很快就会在视觉上变得复杂。上一节的非参数结构方程没有这一缺点，因此，在某些情形下可能更为合适。此外，因果图并不针对某个特定的感兴趣目标参数，因而，由因果图得到的可识别性假设可能强于实际所需。

我们先在图 2.1 中给出式 (2.1) 的一种可能因果图。通过为每个 $X_j$ 定义 $\operatorname{Pa}(X_j)$，并定义联合分布 $P_U$，我们作出因果假设。根据 $f$ 所给出的关系，从 $\operatorname{Pa}(X_j)$ 中的每个节点向 $X_j$ 画有向箭头，并从每个 $U_{X_j}$ 向对应的 $X_j$ 画有向箭头。理想情况下，这些假设应以我们关于感兴趣科学问题的领域知识为依据。直接效应（direct effect），例如 $L_0$ 对 $L_1$ 的直接效应，由两个节点之间的有向箭头表示。如果不确定两个变量之间是否存在直接效应，我们默认将其包含在内，因为明确省略某个箭头，就意味着已知该直接效应不存在。例如，图 2.2 给出了 SCM (2.2) 的一种可能因果图。根据关于 SCM 中随机化的先验知识，我们明确省略了 $L_0$ 与 $A_0$ 之间以及 $L_1$ 与 $A_1$ 之间的有向箭头。

\begin{figure}[htbp]
 \centering
 % 图 2.1；依据原书第 19 页的实线箭头逐条重绘。
\begin{tikzpicture}[causal,x=1.7cm,y=1.3cm,
 ch02n/.style={inner sep=2pt,fill=white,font=\normalsize}]
\node[ch02n] (ch02aL0) at (0,0) {$L_0$};
\node[ch02n] (ch02aA0) at (1,1) {$A_0$};
\node[ch02n] (ch02aL1) at (2,0) {$L_1$};
\node[ch02n] (ch02aA1) at (3,1) {$A_1$};
\node[ch02n] (ch02aY) at (4,0) {$Y$};
\node[ch02n] (ch02aUL0) at (0,-1) {$U_{L_0}$};
\node[ch02n] (ch02aUA0) at (1,2) {$U_{A_0}$};
\node[ch02n] (ch02aUL1) at (2,-1) {$U_{L_1}$};
\node[ch02n] (ch02aUA1) at (3,2) {$U_{A_1}$};
\node[ch02n] (ch02aUY) at (4,-1) {$U_Y$};
\begin{scope}[on background layer]
\draw[ed] (ch02aUL0) -- (ch02aL0);
\draw[ed] (ch02aUA0) -- (ch02aA0);
\draw[ed] (ch02aUL1) -- (ch02aL1);
\draw[ed] (ch02aUA1) -- (ch02aA1);
\draw[ed] (ch02aUY) -- (ch02aY);
\draw[ed] (ch02aL0) -- (ch02aA0);
\draw[ed] (ch02aL0) -- (ch02aL1);
\draw[ed] (ch02aA0) -- (ch02aL1);
\draw[ed] (ch02aA0) -- (ch02aA1);
\draw[ed] (ch02aA0) -- (ch02aY);
\draw[ed] (ch02aL1) -- (ch02aA1);
\draw[ed] (ch02aL1) -- (ch02aY);
\draw[ed] (ch02aA1) -- (ch02aY);
\end{scope}
\end{tikzpicture}

 \caption{式 (2.1) 的一种可能因果图}
 \label{fig:ch02-causal-21}
\end{figure}

\begin{figure}[htbp]
 \centering
 % 图 2.2；保留源图 U_{A_1} -> A_1，不据式 (2.2) 擅改节点。
\begin{tikzpicture}[causal,x=1.7cm,y=1.3cm,
 ch02n/.style={inner sep=2pt,fill=white,font=\normalsize}]
\node[ch02n] (ch02bL0) at (0,0) {$L_0$};
\node[ch02n] (ch02bA0) at (1,1) {$A_0$};
\node[ch02n] (ch02bL1) at (2,0) {$L_1$};
\node[ch02n] (ch02bA1) at (3,1) {$A_1$};
\node[ch02n] (ch02bY) at (4,0) {$Y$};
\node[ch02n] (ch02bUL0) at (0,-1) {$U_{L_0}$};
\node[ch02n] (ch02bUA0) at (1,2) {$U_{A_0}$};
\node[ch02n] (ch02bUL1) at (2,-1) {$U_{L_1}$};
\node[ch02n] (ch02bUA1) at (3,2) {$U_{A_1}$};
\node[ch02n] (ch02bUY) at (4,-1) {$U_Y$};
\begin{scope}[on background layer]
\draw[ed] (ch02bUL0) -- (ch02bL0);
\draw[ed] (ch02bUA0) -- (ch02bA0);
\draw[ed] (ch02bUL1) -- (ch02bL1);
\draw[ed] (ch02bUA1) -- (ch02bA1);
\draw[ed] (ch02bUY) -- (ch02bY);
\draw[ed] (ch02bL0) -- (ch02bL1);
\draw[ed] (ch02bA0) -- (ch02bL1);
\draw[ed] (ch02bA0) -- (ch02bA1);
\draw[ed] (ch02bA0) -- (ch02bY);
\draw[ed] (ch02bL1) -- (ch02bY);
\draw[ed] (ch02bA1) -- (ch02bY);
\end{scope}
\end{tikzpicture}

 \caption{式 (2.2) 的一种可能因果图}
 \label{fig:ch02-causal-22}
\end{figure}

% 本段源书把前文定义为外生的 U 称为 endogenous；此处保留原文用词，详见校对报告。
图 2.1 和图 2.2 都没有在内生误差 $U=(U_{L_0},U_{A_0},U_{L_1},U_{A_1},U_Y)$ 之间画出任何箭头。这表明，式 (2.1) 和式 (2.2) 已经编码了一个关于内生误差联合独立性的强假设。然而，实践中我们不太可能作出这一假设。当希望反映内生变量 $U$ 之间的关系时，我们用虚线双向箭头来表示。如果对式 (2.1) 中 $P_U$ 的分布不作任何假设，所得因果图就如图 2.3 所示。该图中，处理节点与结局 $Y$ 之间有许多所谓的后门路径（backdoor path）。为了分离出感兴趣的因果效应，我们必须阻断从 $A_{0:1}$ 到 $Y$ 的所有尚未被阻断的后门路径。

\begin{figure}[htbp]
 \centering
 % 图 2.3；与图 2.1 相同的 13 条实线有向边，及五个 U 节点间的 10 条虚线双向边。
\begin{tikzpicture}[causal,x=1.7cm,y=1.3cm,
 ch02n/.style={inner sep=2.2pt,fill=white,font=\normalsize},
 ch02dep/.style={<->,line width=.45pt,dash pattern=on 4pt off 4pt}]
\node[ch02n] (ch02cL0) at (0,0) {$L_0$};
\node[ch02n] (ch02cA0) at (1,1) {$A_0$};
\node[ch02n] (ch02cL1) at (2,0) {$L_1$};
\node[ch02n] (ch02cA1) at (3,1) {$A_1$};
\node[ch02n] (ch02cY) at (4,0) {$Y$};
\node[ch02n] (ch02cUL0) at (0,-1) {$U_{L_0}$};
\node[ch02n] (ch02cUA0) at (1,2) {$U_{A_0}$};
\node[ch02n] (ch02cUL1) at (2,-2) {$U_{L_1}$};
\node[ch02n] (ch02cUA1) at (3,2) {$U_{A_1}$};
\node[ch02n] (ch02cUY) at (4,-1) {$U_Y$};
\begin{scope}[on background layer]
\draw[ed] (ch02cUL0) -- (ch02cL0);
\draw[ed] (ch02cUA0) -- (ch02cA0);
\draw[ed] (ch02cUL1) -- (ch02cL1);
\draw[ed] (ch02cUA1) -- (ch02cA1);
\draw[ed] (ch02cUY) -- (ch02cY);
\draw[ed] (ch02cL0) -- (ch02cA0);
\draw[ed] (ch02cL0) -- (ch02cL1);
\draw[ed] (ch02cA0) -- (ch02cL1);
\draw[ed] (ch02cA0) -- (ch02cA1);
\draw[ed] (ch02cA0) -- (ch02cY);
\draw[ed] (ch02cL1) -- (ch02cA1);
\draw[ed] (ch02cL1) -- (ch02cY);
\draw[ed] (ch02cA1) -- (ch02cY);
% 每一对误差节点恰有一条双向边；曲率仅为排版调整。
\draw[ch02dep] (ch02cUA0.north east)
 .. controls (1.5,2.4) and (2.5,2.4) .. (ch02cUA1.north west);
\draw[ch02dep] (ch02cUA0.west)
 .. controls (-1.1,2.15) and (-1.35,.1) .. (ch02cUL0.west);
\draw[ch02dep] (ch02cUA0.west)
 .. controls (-2.2,.65) and (-1.8,-2.8) .. (ch02cUL1.west);
\draw[ch02dep] (ch02cUA0.north east)
 .. controls (5.3,4.7) and (5.55,.3) .. (ch02cUY.north east);
\draw[ch02dep] (ch02cUA1.north west)
 .. controls (-1.3,4.7) and (-1.55,.3) .. (ch02cUL0.north west);
\draw[ch02dep] (ch02cUA1.east)
 .. controls (6.2,.65) and (5.8,-2.8) .. (ch02cUL1.east);
\draw[ch02dep] (ch02cUA1.east)
 .. controls (5.1,2.15) and (5.35,.1) .. (ch02cUY.east);
\draw[ch02dep] (ch02cUL0.south east)
 .. controls (.55,-1.65) and (1.35,-2.1) .. (ch02cUL1.west);
\draw[ch02dep] (ch02cUL1.east)
 .. controls (2.65,-2.1) and (3.45,-1.65) .. (ch02cUY.south west);
\draw[ch02dep] (ch02cUL0.south east)
 .. controls (1,-2.95) and (3,-2.95) .. (ch02cUY.south west);
\end{scope}
\end{tikzpicture}

 \caption{对 $P_U$ 的分布不作任何假设时，式 (2.1) 的因果图}
 \label{fig:ch02-causal-23}
\end{figure}
\FloatBarrier
